\section{Funtores «Tor» locales y globales; fórmula de Künneth}
\label{section:III.6}

\subsection{Introducción}
\label{subsection:III.6.1}

\oldpage[III]{137}

\begin{env}[6.1.1]
\label{III.6.1.1}
Sean $f:X\to Y$ un morfismo de preesquemas y $\mathscr F$ un
$\mathscr O_X$-módulo cuasicoherente. En el estudio de las «imágenes
directas superiores» $\mathrm R^n f_*(\mathscr F)$, se llega a considerar el
siguiente problema general: dado un morfismo de «cambio de base»
$g:Y'\to Y$, se pone
\[
  X'=X_{(Y')}=X\times_Y Y',\qquad
  \mathscr F'=\mathscr F\otimes_Y\mathscr O_{Y'},\qquad
  f'=f_{(Y')}:X'\to Y',
\]
y se pretende obtener información sobre las imágenes directas superiores
$\mathrm R^n f'_*(\mathscr F')$ (suponiendo conocidas las
$\mathrm R^n f_*(\mathscr F)$). Se ve fácilmente (cf., por ejemplo, (7.7.2))
que el problema se reduce a estudiar las variaciones de
$\mathrm R^n f_*(\mathscr F\otimes_{\mathscr O_Y}\mathscr G)$ para un
$\mathscr O_Y$-módulo cuasicoherente variable $\mathscr G$, es decir, el
funtor
\[
  \mathscr G\longmapsto
  \mathrm R^n f_*(\mathscr F\otimes_{\mathscr O_Y}\mathscr G).
\]
Si $\mathscr F$ es plano sobre $Y$, el funtor
$\mathscr G\mapsto\mathscr F\otimes_{\mathscr O_Y}\mathscr G$ es exacto
(0\textsubscript{I}, 6.7.4) y, por consiguiente, el funtor compuesto
$\mathscr G\mapsto\mathrm R^n f_*(\mathscr F\otimes_{\mathscr O_Y}\mathscr G)$
sigue siendo un funtor cohomológico. Pero esto ya no sucede en el caso
general; para poder aplicar los métodos cohomológicos, se ve uno llevado a
sustituir
$\mathscr G\mapsto\mathrm R^n f_*(\mathscr F\otimes_{\mathscr O_Y}\mathscr G)$
por otros funtores, que esta vez son siempre funtores cohomológicos. Estos
funtores, que generalizan los funtores «Tor» de la teoría de módulos, se
definen en los números 6.3 a 6.7; hay, por lo demás, dos generalizaciones de
este tipo, una «local» y la otra «global», relacionadas por sucesiones
espectrales que se discutirán en el número 6.7; como aplicación de estas
sucesiones espectrales, se obtiene en particular, bajo ciertas condiciones,
una «fórmula de Künneth» que expresa
\[
  \mathrm R^n(f_1\times f_2)_*(\mathscr F_1\otimes_Y\mathscr F_2)
\]
por medio de las imágenes directas superiores
$\mathrm R^p f_{1*}(\mathscr F_1)$ y $\mathrm R^q f_{2*}(\mathscr F_2)$.
Otras sucesiones espectrales (6.8) generalizan las sucesiones espectrales de
asociatividad del funtor «Tor» de módulos; por último, el problema del cambio
de base conduce también a sucesiones espectrales (6.9).
\end{env}

\begin{env}[6.1.2]
\label{III.6.1.2}
Además, se comprueba (6.10) que los funtores cohomológicos
$\mathscr G\mapsto\mathscr T_\bullet(\mathscr G)$ así definidos son
localmente (sobre $Y$) de la forma
\[
  \mathscr G\longmapsto
  \mathscr H_\bullet(\mathscr L_\bullet\otimes_Y\mathscr G),
\]
donde $\mathscr L_\bullet$ es un complejo de $\mathscr O_Y$-módulos
localmente libres (definido salvo homotopía) y $\mathscr H_\bullet$ denota la
homología. Conviene entonces dejar de lado la situación particular que dio
origen a $\mathscr T_\bullet$ y estudiar, de

\oldpage[III]{138}

manera general, los funtores de la forma precedente
$\mathscr G\mapsto\mathscr H_\bullet(\mathscr L_\bullet\otimes_Y\mathscr G)$
(imponiendo además, cuando proceda, hipótesis de finitud adecuadas sobre los
$\mathscr L_i$ o los $\mathscr H_i(\mathscr L_\bullet)$): este es el objeto
del § 7, cuya lectura es, en lo esencial, independiente del § 6. Las
propiedades más importantes de estos funtores conciernen a las propiedades de
exactitud de una componente $\mathscr T_i$ de $\mathscr T_\bullet$; se darán
diversos criterios que permiten establecer tales propiedades; como
aplicación, se obtendrán condiciones que permiten afirmar (con las notaciones
de (6.1.1)) que el funtor
$\mathscr G\mapsto\mathrm R^n f_*(\mathscr F\otimes_{\mathscr O_Y}\mathscr G)$
es exacto (lo que se expresará diciendo que $\mathscr F$ es
cohomológicamente plano sobre $Y$ en dimensión $n$). Otra propiedad
importante de las componentes $\mathscr T_i$ de
$\mathscr T_\bullet(\mathscr G)=
\mathscr H_\bullet(\mathscr L_\bullet\otimes_Y\mathscr G)$ es una propiedad
de semicontinuidad de la función
$y\mapsto\dim_{k(y)}(T_i(k(y)))$; cuando $\mathscr T_i$ es exacto, esta
propiedad es sustituida por una propiedad de continuidad, y el recíproco es
además cierto, por Grauert, cuando $Y$ es reducido (7.8.4).
\end{env}

\begin{env}[6.1.3]
\label{III.6.1.3}
En los §§ 6 y 7 hemos utilizado sistemáticamente la hipercohomología, tomando
en todas partes como argumentos complejos de haces en lugar de haces, aunque
la necesidad de este punto de vista solo se manifieste en capítulos
posteriores. El formalismo cohomológico desarrollado con este motivo se hará,
por lo demás, más transparente en el capítulo de este Tratado que estará
dedicado al desarrollo de un álgebra de funtores cohomológicos de haces
coherentes, incluido el formalismo de la dualidad. Pero esto requerirá
desarrollos que exceden el marco del presente capítulo.
\end{env}

\begin{env}[6.1.4]
\label{III.6.1.4}
Para abreviar, dados dos complejos $K^\bullet$, $K'^\bullet$ en una categoría
abeliana $C$, diremos que un morfismo de complejos
$f:K^\bullet\to K'^\bullet$ es una equivalencia homotópica si existe un
morfismo $g:K'^\bullet\to K^\bullet$ tal que los morfismos compuestos
$f\circ g$ y $g\circ f$ sean ambos homotópicos a la identidad (por abuso de
lenguaje, cuando exista una equivalencia homotópica de este tipo, diremos
asimismo que $K^\bullet$ y $K'^\bullet$ son homotópicamente equivalentes).
Cuando se puede definir la hipercohomología de un funtor covariante aditivo
$T$ de $C$ en una categoría abeliana $C'$, respecto de un complejo de $C$
(0, 11.4.3), es inmediato que una equivalencia homotópica
$K^\bullet\to K'^\bullet$ de complejos de $C$ define canónicamente un
isomorfismo
\[
  \mathrm R^\bullet T(K^\bullet)\xrightarrow{\sim}
  \mathrm R^\bullet T(K'^\bullet)
\]
en hipercohomología (loc. cit.).
\end{env}
